\problema{Course Schedule}{207}{src/01-grafos/207-course-schedule.cpp}{M}

Hay \texttt{numCourses} cursos numerados de \texttt{0} a
\texttt{numCourses - 1}, y una lista \texttt{prerequisites} de pares
$[a, b]$ que significan ``para tomar el curso $a$ primero hay que tomar el
curso $b$''. Queremos saber si es posible tomar todos los cursos.

\paragraph{La idea, con un ejemplo de la vida real.} Piensa en los cursos como
en armar una torre de bloques de juguete. Para poner un bloque arriba, primero
tiene que estar puesto el bloque de abajo. Cada par $[a, b]$ quiere decir
``antes del bloque $a$ hay que poner el bloque $b$''. Si dibujas una flecha de
$b$ hacia $a$ (\emph{$b$ va antes que $a$}) para cada regla, terminas con un
dibujo de puntos y flechas. En matemáticas ese dibujo se llama \emph{grafo
dirigido}: los puntos son los cursos y las flechas dicen quién va antes que
quién.

Ahora, la pregunta ``¿puedo tomar todos los cursos?'' es lo mismo que
preguntar ``¿este dibujo tiene un círculo vicioso?''. Un \emph{círculo vicioso}
(en matemáticas, un \emph{ciclo}) es cuando el curso A pide al B, el B pide al
C, y el C vuelve a pedir al A: nadie puede ser el primero, así que es imposible.
Si no hay ningún círculo vicioso, siempre puedes acomodar los cursos en una
fila donde cada uno queda después de todos los que necesita. A esa fila
ordenada se le llama \emph{orden topológico}.

\begin{center}
\ttfamily
\begin{tabular}{c}
Con ciclo (imposible):\\[4pt]
$0 \rightarrow 1 \rightarrow 0$ \quad (0 necesita 1, 1 necesita 0)
\end{tabular}
\qquad
\begin{tabular}{c}
Sin ciclo (posible):\\[4pt]
$0 \rightarrow 1 \rightarrow 3$,\ \ $0 \rightarrow 2 \rightarrow 3$\\
orden topológico: $0, 1, 2, 3$
\end{tabular}
\end{center}

\paragraph{Una palabra nueva: ``grado de entrada''.} El \emph{grado de entrada}
de un curso es simplemente cuántos prerrequisitos le faltan todavía por
cumplir, es decir, cuántas flechas le llegan apuntando. Un curso con grado de
entrada $0$ no le debe nada a nadie, así que lo puedes tomar de inmediato.
Piénsalo como una lista de tareas: las que no dependen de ninguna otra son las
que puedes empezar hoy mismo.

\paragraph{Cómo lo resuelve el código, paso a paso.} La técnica se llama
\emph{algoritmo de Kahn}, pero la idea es muy sencilla: ir haciendo primero los
cursos que ya no necesitan nada. Usamos una \emph{cola}, que es como la fila
del súper: el primero que entra es el primero que sale.
\begin{enumerate}
  \item Anotamos, para cada curso, la lista de cursos que dependen de él
        (\texttt{adj[b]} guarda los cursos que necesitan a \texttt{b}) y su
        grado de entrada (cuántos prerrequisitos le faltan).
  \item Metemos en la cola todos los cursos que ya tienen grado de entrada $0$,
        porque esos se pueden tomar sin esperar a nadie.
  \item Sacamos un curso de la cola, lo contamos como ``ya tomado'' y a cada
        curso que dependía de él le restamos $1$ prerrequisito. Si a alguno se
        le acaban los prerrequisitos pendientes (llega a $0$), lo metemos a la
        cola.
  \item Repetimos hasta vaciar la cola.
\end{enumerate}
Al final contamos: si logramos tomar todos los cursos, no había ningún círculo
vicioso, así que la respuesta es sí. Si sobró algún curso (nunca bajó su grado
de entrada a $0$ porque estaba atrapado en un círculo), la respuesta es no.

\lstinputlisting{src/01-grafos/207-course-schedule.cpp}

\complejidad{$O(V+E)$}{$O(V+E)$}

\paragraph{¿Qué significa esa complejidad?} Aquí $V$ es el número de cursos (los
puntos del dibujo) y $E$ el número de reglas de prerrequisito (las flechas).
Decir que el tiempo es $O(V+E)$ solo quiere decir que miramos cada curso una
vez y cada flecha una vez, y ya. No repetimos trabajo, así que es lo más rápido
que se puede: para saber si el plan cabe, hay que leer al menos todas las reglas
una vez.

\paragraph{Consejos para la entrevista (Amazon).} Este problema es la puerta de
entrada al patrón de \emph{orden topológico}, que aparece mucho cuando hay que
planear tareas que dependen unas de otras (por ejemplo, en qué orden compilar
las partes de un programa o instalar programas que necesitan otros programas).
Menciona que se puede resolver de dos maneras: contando prerrequisitos con una
cola (el algoritmo de Kahn, el que usamos aquí, que además te sirve si te piden
el orden concreto), o con \emph{DFS} (recorrer hacia el fondo por las flechas)
pintando los cursos de colores para cachar el círculo vicioso. No olvides los
casos raros: sin prerrequisitos (siempre se puede), un círculo vicioso que
atrapa solo a algunos cursos y no a todos, y cursos que tienen muchos
prerrequisitos o de los que dependen muchos otros.
